\subsection{张量积上的算子；张量积作为多重模}

在第1号的假设和记号下，对于每个自同态 \(u\) （相应地 \(v\) ，即 A-模 \(E\) （相应地 \(F\) ) 的自同态）， \(u \otimes l_F\) （相应地 \(l_E \otimes v\) ）是 \(Z\) -模 \(E \otimes_A F\) 的自同态；由第2号的公式 (5) 立即得出，映射 \(u \mapsto u \otimes l_F\) （相应地 \(v \mapsto l_E \otimes v\) ）是一个环同态

\[
\operatorname{End} _ {\mathbf {A}} (\mathbf {E}) \rightarrow \operatorname{End} _ {\mathbf {Z}} (\mathbf {E} \otimes_ {\mathbf {A}} \mathbf {F})
\]

（相应地 \(\operatorname{End}_{\mathbf{A}}(\mathbf{F}) \to \operatorname{End}_{\mathbf{Z}}(\mathbf{E} \otimes_{\mathbf{A}} \mathbf{F}))\) ；此外，

\[
(u \otimes \mathrm{l} _ {\mathrm{F}}) \circ (\mathrm{l} _ {\mathrm{E}} \otimes v) = (\mathrm{l} _ {\mathrm{E}} \otimes v) \circ (u \otimes \mathrm{l} _ {\mathrm{F}}) = u \otimes v\tag{7}
\]

因此（§ 1，第 14 号） E ⊗\_A F 关于环 End\_A(E) 与 End\_A(F) 具有典范的双模结构。

既然如此，假设在 E 上给定一个 \(((\mathbf{B}_i^{\prime});\mathbf{A},(\mathbf{C}_j^{\prime}))\) -多重模结构，在 F 上给定一个 \((\mathbf{A},(\mathbf{B}_h^{\prime});(\mathbf{C}_k^{\prime}))\) -多重模结构（§ 1，第 14 号）；这等同于说，给定了环同态 \(\mathrm{B}_i^\prime \to \mathrm{End}_{\mathbf{A}}(\mathbf{E})\) , \(\mathrm{C}_j^{\prime 0} \to \mathrm{End}_{\mathbf{A}}(\mathbf{E})\) ，其像两两可交换，以及环同态 \(\mathrm{B}_h^{\prime \prime} \to \mathrm{End}_{\mathbf{A}}(\mathbf{F})\) , \(\mathrm{C}_k^{\prime \prime 0} \to \mathrm{End}_{\mathbf{A}}(\mathbf{F})\) ，其像也两两可交换。若将这些同态分别与上面定义的典范同态 \(\mathrm{End}_{\mathbf{A}}(\mathbf{E}) \to \mathrm{End}_{\mathbf{Z}}(\mathbf{E} \otimes_{\mathbf{A}} \mathbf{F})\) 和 \(\mathrm{End}_{\mathbf{A}}(\mathbf{F}) \to \mathrm{End}_{\mathbf{Z}}(\mathbf{E} \otimes_{\mathbf{A}} \mathbf{F})\) 复合，则（考虑到 (7)）可以看出，环同态

\[
\mathrm{B} _ {i} ^ {\prime} \rightarrow \operatorname{End} _ {\mathbf {Z}} (\mathrm{E} \otimes_ {\mathrm{A}} \mathrm{F}), \quad \mathrm{C} _ {j} ^ {\prime 0} \rightarrow \operatorname{End} _ {\mathbf {Z}} (\mathrm{E} \otimes_ {\mathrm{A}} \mathrm{F})
\]

\[
\mathrm{B} _ {h} ^ {\prime \prime} \rightarrow \operatorname{End} _ {\mathbf {Z}} (\mathrm{E} \otimes_ {\mathrm{A}} \mathrm{F}), \quad \mathrm{C} _ {k} ^ {\prime \prime 0} \rightarrow \operatorname{End} _ {\mathbf {Z}} (\mathrm{E} \otimes_ {\mathrm{A}} \mathrm{F})
\]

其像两两可交换；换句话说，在 \(\mathbf{E} \otimes_{\mathbf{A}} \mathbf{F}\) 上定义了一个 \(((\mathbf{B}_i'), (\mathbf{B}_h''); (\mathbf{C}_j'), (\mathbf{C}_k'')\) -多重模结构；正是这个多重模也称为 \(((\mathbf{B}_i^{\prime});\mathbf{A},(\mathbf{C}_j^{\prime}))\) -多重模 E 与 \((\mathbf{A},(\mathbf{B}_h^{\prime\prime});(\mathbf{C}_k^{\prime\prime}))\) -多重模 F 的（相对于 A 的）张量积。这个多重模是类似于第1号中所考虑的泛映射问题的解；确切地说：

命题3. 设 G 为 \(((\mathbf{B}_i'), (\mathbf{B}_h''); (\mathbf{C}_j'), (\mathbf{C}_k''))-multimodule.\) 。

\begin{enumerate}
\def\labelenumi{(\alph{enumi})}
\tightlist
\item
  设 \(g\) 是多重模 \(\mathbf{E} \otimes_{\mathbf{A}} \mathbf{F}\) 到 \(\mathbf{G}\) 的一个线性映射。映射 \(f: (x, y) \mapsto g(x \otimes y)\) （从 \(\mathbf{E} \times \mathbf{F}\) 到 \(\mathbf{G}\) ）是 \(\mathbf{Z}\) -双线性的，且满足第1号的关系 (1) 以及条件
\end{enumerate}

\[
\left\{ \begin{array}{l l} f (\mu_ {i} ^ {\prime} x, y) = \mu_ {i} ^ {\prime} f (x, y), & f (x \nu_ {j} ^ {\prime}, y) = f (x, y) \nu_ {j} ^ {\prime} \\ f (x, \mu_ {h} ^ {\prime \prime} y) = \mu_ {h} ^ {\prime \prime} f (x, y), & f (x, y \nu_ {k} ^ {\prime \prime}) = f (x, y) \nu_ {k} ^ {\prime \prime} \end{array} \right.\tag{8}
\]

对所有 \(x \in \mathbf{E}, y \in \mathbf{F}\) , \(\mu_i' \in \mathbf{B}_i'\) , \(\nu_j' \in \mathbf{C}_j'\) , \(\mu_h'' \in \mathbf{B}_h''\) , \(\nu_k'' \in \mathbf{C}_{k}''\) 成立，其中 \(i, j, h, k\) 任意。

\begin{enumerate}
\def\labelenumi{(\alph{enumi})}
\setcounter{enumi}{1}
\tightlist
\item
  反之，设 \(f\) 是一个 \(\mathbf{Z}\) -双线性映射（从 \(\mathbf{E} \times \mathbf{F}\) 到 \(\mathbf{G}\) ），满足条件 (1)（第1号）和 (8)。则存在且仅存在一个线性映射 \(g\) （从多重模 \(\mathbf{E} \otimes_{\mathbf{A}} \mathbf{F}\) 到多重模 \(\mathbf{G}\) ），使得 \(f(x, y) = g(x \otimes y)\) ，其中 \(x \in \mathbf{E}, y \in \mathbf{F}\) 。
\end{enumerate}

断言 (a) 直接由 \(\mathbf{E} \otimes_{\mathbf{A}} \mathbf{F}\) 上多重模结构的定义得出，因为例如 \((x \otimes y)\nu_j' = (x\nu_j') \otimes y\) 。为证明 (b)，我们首先注意到，第1号的命题1给出了 \(\mathbf{Z}\) -线性映射 \(g\) 的存在性和唯一性，它满足 \(g(x \otimes y) = f(x, y)\) （其中 \(x \in \mathbf{E}\) , \(y \in \mathbf{F}\) ）；只需验证 \(g\) 关于多重模结构是线性的。由于元素 \(x \otimes y\) 生成 \(\mathbf{Z}\) -模 \(\mathbf{E} \otimes_{\mathbf{A}} \mathbf{F}\) ，只需验证关系 \(g(\mu'(x \otimes y)) = \mu_i' g(x \otimes y)\) 及其类似关系；但这由公式 \(g(x \otimes y) = f(x, y)\) 以及关系式 (8) 立即得出。

注释。\(\mathbf{E} \otimes_{\mathbf{A}} \mathbf{F}\) 的元素一般说来可以有多种方式写成 \(\sum_{i} (x_i \otimes y_i)\) 的形式，其中 \(x_i \in \mathbf{E}\) 且 \(y_i \in \mathbf{F}\) ；但要定义线性映射 \(g\) （从多重模 \(\mathbf{E} \otimes_{\mathbf{A}} \mathbf{F}\) 到一个多重模 \(\mathbf{G}\) ），无需验证：若 \(\sum_{i} (x_i \otimes y_i) = \sum_{j} (x_j' \otimes y_j')\) ，则 \(\sum_{i} g(x_i \otimes y_i) = \sum_{j} g(x_j' \otimes y_j')\) ；只需给定 \(g(x \otimes y)\) （对于 \(x \in \mathbf{E}\) 与 \(y \in \mathbf{F}\) ），并验证 \((x, y) \mapsto g(x \otimes y)\) 是 \(\mathbf{Z}\) -双线性的且满足条件 (1)（第1号）和 (8)。

设 \(\mathbf{E}'\) 是 \(((\mathbf{B}_i'),\mathbf{A},(\mathbf{C}_j')\) -多重模， \(\mathbf{F}'\) 一个 \((\mathbf{A},(\mathbf{B}_h'');(\mathbf{C}_k'')\) -多重模，且 \(u:\mathbf{E}\to \mathbf{E}',v:\mathbf{F}\to \mathbf{F}'\) 为多重模的线性映射；由定义（第2号）立即得出， \(u\otimes v\) 是从多重模 \(\mathbf{E}\otimes_{\mathbf{A}}\mathbf{F}\) 到多重模 \(\mathbf{E}'\otimes_{\mathbf{A}}\mathbf{F}'\) 的线性映射。

设 E 始终表示右 A-模，令 \(_{s}A_{d}\) 表示把环 A 视为 (A, A)-双模（ \(\S\) 1，第 14 号，例 1）；由上述可知，张量积 \(\mathbf{E} \otimes_{\mathbf{A}} (\mathbf{s}\mathbf{A}_{d})\) 具有典范的右 A-模结构，使得 \((x \otimes \lambda)\mu = x \otimes (\lambda\mu)\) ，其中 \(x \in E\) , \(\lambda \in A\) , \(\mu \in A\) 。映射 \((x, \lambda) \mapsto x\lambda\) （从 \(\mathbf{E} \times (\mathbf{s}\mathbf{A}_{d})\) 到 E ）是 Z-双线性的，并满足条件 (1)（第 1 号）和 (8)

（在后一种情形中， \(\mathbf{B}_i'\) , \(\mathbf{C}_j'\) 和 \(\mathbf{B}_h''\) 不出现，而族 \((\mathbf{C}_k'')\) 就化为 \(\mathbf{A}\) ）；因此（命题3），存在一个 A-线性映射 \(g\) （称为典范映射），它把 \(\mathbf{E} \otimes_{\mathbf{A}} (s\mathbf{A}_d)\) 映入 \(\mathbf{E}\) ，使得 \(g(x \otimes \lambda) = x\lambda\) ，其中 \(x \in \mathbf{E}\) , \(\lambda \in \mathbf{A}\) 。

命题4. 若 E 是一个右 A-模，则映射 \(h: x \mapsto x \otimes 1\) （它把 E 映入 \(\mathbf{E} \otimes_{\mathbf{A}} (s\mathbf{A}_d)\) ）是一个右 A-模同构，其逆同构 \(g\) 满足 \(g(x \otimes \lambda) = x\lambda\) ，其中 \(x \in \mathbf{E}, \lambda \in \mathbf{A}\) 。

若 \(g\) 是典范映射，则 \(g \circ h\) 是恒等映射 \(1_{\mathbf{E}}\) ，且 \(h \circ g\) 与 \(\mathbf{E} \otimes_{\mathbf{A}} (s\mathbf{A}_d)\) 到自身的恒等映射在形如 \(x \otimes y\) 的元素上一致，而这些元素生成这个 \(\mathbf{Z}\) -模；由此得出结论。

我们通常将写成 \(E \otimes_{A} A\) 而不是 \(E \otimes_{A} (s A_{d})\) ，并经常通过上述典范同构把 \(E \otimes_{A} A\) 与 E 等同。注意，若 E 还具有与其右 A-模结构相容的（左或右）B-模结构，则 g 和 h 对于 E 和 \(E \otimes_{A} A\) 上的 B-模结构也是同构（从而也是多重模同构）。

现在设 F 是一个左 A-模； \((sA_{d}) \otimes_{A} F\) （也记作 A \(\otimes_{A} F\) ）则具有典范的左 A-模结构，并且如同命题 4 那样，可定义一个从 A \(\otimes_{A} F\) 到 F 上的典范同构，它把 \(\lambda \otimes x\) 映为 \(\lambda x\) ，其逆同构为 \(x \mapsto 1 \otimes x\) 。

特别地，存在 (A, A)-双模 \((_{s}\mathbf{A}_{d})\otimes_{\mathbf{A}}(_{s}\mathbf{A}_{d})\) 到 \({}_{s}\mathbf{A}_{d}\) 上的一个典范同构，它把 \(\lambda \otimes \mu\) 映为 \(\lambda \mu\) 。

